% Autore: Federico Gandellini
% Titolo: mole.io

% Documento
\documentclass[a4paper,11pt,italian]{report}
%\documentclass[a4paper,11pt,italian,draft]{report} % Per non caricare le immagini

% Packages
\usepackage[italian]{babel}
\usepackage[latin1]{inputenc}
\usepackage{graphicx}
\usepackage{color}
\usepackage{colortbl}
\usepackage{amsmath}
\usepackage{amsfonts}
\usepackage{rotating}

%% Formato proposto dall'universit�
%\documentclass[12pt,fleqn,twoside,a4paper]{book}
%\usepackage[latin1]{inputenc} %% lettere accentate
%\usepackage[italian]{babel}
%\usepackage[top=2.5cm,margin=3cm,bottom=2.5cm]{geometry}
%\usepackage{graphicx}
%\usepackage{palatino}
%\frontmatter
%\makeindex
%\mainmatter

% Numeri di pagina e intestazione
\usepackage{fancyhdr}
% Line spacing -----------------------------------------------------------
\newlength{\defbaselineskip}
\setlength{\defbaselineskip}{\baselineskip}
\newcommand{\setlinespacing}[1]%
           {\setlength{\baselineskip}{#1 \defbaselineskip}}
\newcommand{\doublespacing}{\setlength{\baselineskip}%
                           {2.0 \defbaselineskip}}
\newcommand{\singlespacing}{\setlength{\baselineskip}{\defbaselineskip}}

%-----------------------------FANCYHDR--------------------------------------%
\pagestyle{fancy}                       % Sets fancy header and footer
\fancyfoot{}                            % Delete current footer settings
\renewcommand{\chaptermark}[1]{         % Lower Case Chapter marker style
  \markboth{\chaptername\ \thechapter.\ #1}{}} %
\renewcommand{\sectionmark}[1]{         % Lower case Section marker style
  \markright{\thesection.\ #1}}         %
\fancyhead[LE,RO]{\bfseries\thepage}    % Page number (boldface) in left on even
                                        % pages and right on odd pages
\fancyhead[RE]{\bfseries\leftmark}      % Chapter in the right on even pages
\fancyhead[LO]{\bfseries\rightmark}     % Section in the left on odd pages
\renewcommand{\headrulewidth}{0.05pt}    % Width of head rule
%\pagestyle{fancy}\addtolenght{\headwidth}{20pt}
%\rnewcommand{\chaptermark}[1]{\markboth{\thechapter.\#1}{}}
%\rnewcommand{\sectionmark}[1]{\markright{\thesection\#1}{}}
%\cfoot{}
%\rhead[\fancyplain{}{\bfseries\leftmark}]{\fancyplain{}{\bfseries\thepage}}
%\lhead[\fancyplain{}{\bfseries\thepage}]{\fancyplain{}{\bfseries\rightmark}}


% Formattazione
\linespread{1.6} % Interlinea 2

\begin{document}

\begin{titlepage}
  \begin{center}
    \includegraphics[height=5.0cm]{img/minerva_2013_DI.jpg}
    
    \vspace*{.4cm}
    {\Large 
      \emph{Corso di Laurea Magistrale in\\[.3cm]
        Scienze e Tecnologie dell'Informazione}
    }
    \vfill
    \begin{LARGE}
      \textbf{Mole.io\\[0.4cm]
        un Sistema per la Gestione Centralizzata\\[0.6cm]
        dei Log Applicativi}
    \end{LARGE}
    
    \vfill
    \begin{minipage}{.99\linewidth}
      \begin{tabular}{l r}
        \begin{minipage}{.4\linewidth}
          \begin{flushleft}
            {\large
              RELATORE\\[.3cm]
              Prof. Ernesto Damiani
            }
            
            \vspace*{1.5cm}
            {\large
              CORRELATORE\\[.3cm]
              Emanuele DelBono
            }
          \end{flushleft}
        \end{minipage}
        &
        \begin{minipage}{.6\linewidth}
          \begin{flushright}
            {\large
              TESI DI LAUREA DI\\[.3cm]
              Federico Gandellini\\[.45cm]
              Matr. 703156
            }
          \end{flushright}
        \end{minipage}
      \end{tabular}
    \end{minipage}
    
    \vfill
    {\large{{Anno Accademico 2012/2013}}}
  \end{center}
\end{titlepage}

\clearpage
\begin{titlepage}

\setlength{\topmargin}{1pt} \setlength{\headheight}{1pt}

\huge{\textbf{Ringraziamenti}}

\vspace{3.0cm} \normalsize \noindent Un grazie a ...

% parenti e amici
% codice plastico
% chi fa sw open

\end{titlepage}
\clearpage
\tableofcontents
\clearpage

\chapter*{Introduzione\markboth{}{Introduzione}}
\addcontentsline{toc}{chapter}{Introduzione}
In questa tesi si descriver� Mole.io: un sistema centralizzato per la raccolta e l'aggregazione di messaggi provenienti da applicazioni remote.

Durante il loro ciclo di lavoro o \textit{processing}, le applicazioni software eseguono operazioni significative o entrano in situazioni di errore. In questi casi � importante che le persone che hanno in carico la gestione di questi sistemi, siano informate dell'accaduto in modo da operare scelte opportune o applicare le dovute correzioni (\textit{bugfix}).

Gli sviluppatori sono soliti utilizzare messaggi di tracciamento (\textit{log}) per stampare a video o salvare in \textit{file} stati significativi delle applicazioni. I messaggi pi� frequenti riportati nei log sono quelli relativi a situazioni di errore (\textit{Exception} e \textit{Stack Trace}).

L'approccio comune alla creazione e gestione dei log presenta la criticit� specifica della \textit{localit�}, poich� tipicamente questi file vengono salvati nella stessa macchina sulla quale sta operando l'applicazione.

All'aumentare del numero di applicazioni da gestire e del numero di macchine in produzione, capita spesso che i server siano in luoghi geograficamente distanti tra loro. Questa situazione rende evidente la difficolt� di ottenere un \textit{feedback} veloce dello stato di ogni software e delle eventuali situazioni di errore in cui le applicazioni si trovano.

Mole.io cerca di risolvere il problema facendo in modo che i software che lo utilizzano, siano in grado di inviare le informazioni che ritengono significative ad un server centrale, il quale le raccoglie, le cataloga e le aggrega per essere facilmente supervisionate da parte degli sviluppatori.

L'esigenza di una applicazione per la centralizzazione dei log nasce da CodicePlastico \cite{website:CodicePlastico}, una azienda con sede a Brescia, che si occupa di realizzare applicazioni su misura per i propri clienti. 

La gestione di un gran numero installazioni dislocate sul territorio e di molte realt� aziendali con esigenze differenti ha reso, per CodicePlastico, particolarmente complesso il tracciamento dello stato di ogni software in produzione. Questa situazione ha spinto l'azienda a decidere di dotarsi di un sistema centralizzato in grado di collezionare i log prodotti dai diversi applicativi, organizzarli e catalogarli in modo automatico. 

Il nuovo approccio permette agli sviluppatori di identificare, in breve tempo, il manifestarsi di un malfunzionamento in qualunque applicazione installata presso uno dei proprio clienti e reagire rapidamente proponendo una azione risolutiva.

CodicePlastico opera nel settore IT avvalendosi di strumenti software di vario genere, sia proprietari, sia \textit{open source}. Per il \textit{deploy} in produzione di Mole.io, si � scelto di utilizzare \textit{Microsoft Azure}, un sistema \textit{PaaS} distribuito con il quale � possibile creare architetture facilmente scalabili per supportare la variabilit� del carico di lavoro richiesto al sistema. 

mettila gi� cos�: in un fase di test il sistema sar� testato su un numero X di installazione, con in progetto di aumentare il campione successivamente... non siamo umili a priori ;-)

Durante una prima fase di \textit{test}, Mole.io sar� utilizzato per gestire tre o quattro installazioni, che fungeranno da \textit{pilota} per il progetto. La pianificazione dell'azienda, per l'immediato futuro, prevede di aumentare velocemente il numero di clienti attivi nel sistema. Una architettura scalabile, di conseguenza, permetter� di affrontare le esigenze di reattivit� e stabilit� del sistema al variare del carico di lavoro.

Nel primo capitolo saranno trattati approfonditamente la tematica dei log, i contesti nei quali essi vengono utilizzati e le problematiche legate alla gestione di questo tipo di soluzione di tracciamento. Sar� analizzato anche come utilizzare i log per ottenere informazioni di supporto alla \textit{business intelligence}. 

Il secondo capitolo riporter� un elenco dei principali \textit{software} per la gestione centralizzata dei log presenti sul mercato e delle soluzioni \textit{Open Source} che sono state prese a modello per la realizzazione di Mole.io. Verr� descritta ogni applicazione e sar� mostrato come Mole.io possa essere una soluzione innovativa sotto svariati punti di vista.

I due capitoli seguenti permetteranno di approfondire i dettagli tecnici delle metodologie di sviluppo applicate durante il \textit{design} del software e alcune tra le principali tecnologie utilizzate per la realizzazione del sistema.

Il quinto capitolo descriver� la struttura di Mole.io e le varie componenti software che rendono l'applicazione scalabile e garantiscono l'alta accessibilit� della soluzione.

Nel sesto capitolo verr� mostrato in modo oggettivo, con \textit{benchmark} e \textit{stress test} il comportamento di Mole.io all'aumentare del carico di lavoro e sar� dimostrato come le soluzioni di design applicate garantiscano buone \textit{performance}, anche in condizioni critiche di traffico.

Infine verranno discussi i risultati ottenuti e saranno proposte alcune interessanti funzionalit� che trasformeranno Mole.io dall'attuale \textit{proof of concept} ad un vero e proprio servizio.

\clearpage

\chapter{Log:\\contesti e problematiche}\label{Log_contesti_e_problematiche}
Prima di entrare nello specifico delle tematiche trattate, � necessario riprendere e chiarire alcuni concetti chiave che utilizzeremo ripetutamente nel seguito della tesi.

Si definisce \textit{processo} un programma in esecuzione e \textit{log} l'insieme dei messaggi prodotti, a fini informativi, da tale processo.

I messaggi di \textit{log} posso essere di vario tipo. \`{E} usanza comune caratterizzare ogni messaggio con un livello di gravit� (\textit{severity}) permettendo cos� una rapida identificazione degli errori critici allo scopo di rendere repentino l'intervento di riparazione dell'applicazione.

Bench� esista un protocollo riconosciuto a livello internazionale e adottato negli ambienti \textit{Unix-like} chiamato \textit{SysLog}, nell'ambito dello sviluppo di applicativi, non esistono norme vincolanti per la strutturazione dei log stessi. Questo accade sia perch� SysLog non rappresenta uno standard rigidamente definito, sia perch� l'organizzazione delle informazioni contenute nei log, � spesso delegata agli sviluppatori, i quali implementano questa funzionalit� nel modo pi� conveniente rispetto alle specifiche esigenze dell'applicazione in costruzione.  

All'interno dei log, di conseguenza, troveremo informazioni diversificate in base al caso d'uso, ma tra le pi� frequenti possiamo citare:
\begin{itemize}
\item dati specifici del sistema nel quale l'applicazione � in esecuzione;
\item dati relativi all'utente che sta utilizzando l'applicazione;
\item dati relativi allo stato del sistema in un preciso istante temporale;
\item un marcatore temporale (\textit{timestamp})
\item un messaggio in linguaggio naturale, significativo, che lo rende immediatamente identificabile tra altri;
\end{itemize}

	\section{Trattare gli errori applicativi}\label{Trattare_gli_errori_applicativi}
	Il salvataggio dei log, come abbiamo anticipato, � una operazione molto comune nei software, ma diventa fondamentale quando si vuole monitorare lo stato interno di una applicazione in esecuzione, con l'obiettivo di essere informati riguardo alle situazioni di errore nelle quali quest'ultima incorre.

Salvare le informazioni relative alle situazioni di errore � importante per gli sviluppatori. Questa operazione permette, infatti di velocizzare l'individuazione di errori (\textit{bug}) nel flusso di lavoro del programma e, di conseguenza la loro risoluzione (\textit{bugfix}).

Il salvataggio dei \textit{log} avviene tipicamente su uno o pi� \textit{file} di testo presenti nella stessa macchina nella quale sta funzionando l'applicazione. A seconda del tempo di esecuzione di una applicazione e della frequenza con la quale essa produce messaggi di \textit{log}, questi \textit{file} possono diventare molto grandi.

\textit{File} di considerevoli dimensioni sono altamente complessi da gestire da parte degli addetti ai lavori. La problematica pi� evidente diviene infatti trovare informazioni significative all'interno di questa grande mole di dati. Questo processo richiede infatti tempo e attenzione in situazioni di emergenza, nelle quali il ripristino del sistema deve avvenire nel modo pi� rapido possibile.

Il problema dei file di grandi dimensioni non riguarda esclusivamente la scansione sequenziale delle informazioni, bens� l'individuazione del messaggio prodotto a fronte di una criticit� e la correlazione di quest'ultima allo stato del sistema nell'istante in cui � stata generata.

L'individuazione degli errori non riguarda esclusivamente l'analisi dei log nell'istante in cui il malfunzionamento si � manifestato, bens� richiede di comprendere la catena di eventi pregressi, non sempre palese, che ha portato il sistema nella condizione di errore. La capacit� dell'analista sta nel riconoscere pattern ricorrenti che generano l'intera situazione.

Un'ulteriore complicazione dovuta alla dimensione eccessiva dei file di log non riguarda solo il riconoscimento delle cause di un problema ma anche l'individuazione di criticit� simili occorse in istanti temporali differenti. Questo processo, noto come \textit{clustering}, diviene ovviamente oneroso in termini di tempo, al crescere delle dimensione del \textit{log}.

L'analista ha anche un'altra incombenza: verificare che le dimensioni dei file di \textit{log} non eccedano al punto di compromettere il funzionamento dell'applicazione stessa a causa dell'assenza di spazio su disco fisso.

A livello aziendale il tempo che intercorre tra la scoperta di un \textit{bug} e il relativo \textit{bugfix} dovrebbe essere il pi� possibile contenuto. Spesso purtroppo le eccessive dimensioni dei file di \textit{log} rendono questa procedura molto costosa.

 
	\clearpage
	\section{La centralizzazione}\label{La_centralizzazione}
	tante applicazioni (anche tipi diversi) che loggano
tanti clienti da gestire dislocati sul territorio
	\clearpage
	\section{Business intelligence}\label{Businness_Intelligence}
	capire come gli utenti usano il sistema
statistiche sul sistema
decidere come indirizzare lo sviluppo
    \clearpage

\chapter{Log: software e applicazioni}\label{Log_software_e_applicazioni}
Durante la prima fase del progetto si sono analizzate le soluzioni gi� proposte dal mercato IT, verificando quali tra i prodotti esistenti presentassero la peculiarit� della centralizzazione dei \textit{log}.

In questo capitolo analizzeremo alcune tra le principali applicazioni dedicate alla raccolta e analisi dei dati, indicando per ciascuna quali \textit{features} sono state di ispirazione al progetto, cos� come quali problemi e svantaggi sono stati identificati.

Bench� il panorama delle soluzioni esistenti fosse ampio, vedremo nel seguito del capitolo, quali sono state le motivazioni che hanno portato alla scelta dell'implementazione di un software proprietario e interno all'azienda.
	\section{Prodotti e soluzioni sul mercato}\label{Prodotti_e_soluzioni_sul_mercato}
	overview di alcuni sistemi di logging con le relative funzioni specifiche
i competitor
airbreak - logga solo
rollbar - aggrega
papertrail - live log

sul mercato esistono svariati sistemi per la gestione centralizzata dei log


Airbrake
+
aggrega log e ne facilita l'analisi
integrato con sistemi di tracking 
-
formati dei messaggi fisso e solo errori


log.io
+
realtime
-
no aggregazione

rollbar
+
non logga solo errori, ma la gestione delle "non eccezioni" � un o' tirata per i capelli
fa aggregazione

https://rollbar.com/vs/airbrake/ questo � mooooolto meglio di mole :(

-
soluzione completa ma blindata e non estensibile (gniiiiiii)



papertrail
+
realtime
aggregazione
-
logga solo tipi di errori noti, no formato flessibile







CACCHIO! fluentd + elasticsearch + kibana = mole
%(mole � pi� facile....gniiiiii)



%mole � pensato pi� come una piattaforma espandibile, accento sulla possibilit� di aggiungere features a caldo
	\clearpage
	\section{Una nuova applicazione: Mole.io}\label{Una_nuova_applicazione_Moleio}
	L'analisi dello stato dell'arte delle tecnologie presenti sul mercato, ha portato CodicePlastico a perseguire la scelta dello sviluppo di una nuova applicazione realizzata internamente all'azienda stessa.

Le motivazioni che hanno determinato questa scelta, sono molteplici:
\begin{description}
\item[Business] la maggior parte delle soluzioni analizzate � fornita come \textit{Software as a Service} (SaaS). In questo caso sarebbe d'obbligo far adottare tale soluzione ad ogni cliente, perdendo di fatto, il vantaggio di avere un sistema centralizzato di raccolta dei dati.
\item[Riservatezza] le soluzioni analizzate che sono fornite come SaaS, prevedono che i messaggi vengano salvati su basi di dati remote e non di propriet� dell'azienda. Questo espone potenzialmente a problematiche di gestione della riservatezza dei dati sensibili che sarebbero pi� facilmente controllabili utilizzando un sistema \textit{self-hosted}.
\item[Personalizzazione] a parte l'ultima soluzione analizzata (Fluentd, Elastic Search e Kibana) la maggior parte dei software presenta difficolt� nella personalizzazione del contenuto dei messaggi, fondamentale per l'azienda, sia per tracciare con precisione problematiche specifiche delle applicazioni in produzione, sia nell'ottica della raccolta dei dati per la Business Intelligence.
\item[Manutenzione] limitare gli interventi di manutenzione ad un'unica applicazione proprietaria, di cui, conseguentemente, si ha pieno controllo dello sviluppo.
\item[Estensibilit�] l'intero sviluppo delle varie componenti dell'applicazione � stato guidato dal concetto di estensibilit�. Il sistema risultante dovr� essere facilmente estendibile con nuove funzionalit� e \textit{feature} per adattarsi nel modo pi� aderente possibile alle esigenze di ogni cliente. Un ulteriore aspetto che ha motivato alla costruzione di una nuova applicazione, � la possibilit� di rilasciare gradualmente agli utenti le nuove funzionalit� introdotte, man mano, nel sistema (\textit{deploy} graduale).
\item[Sperimentazione] tra i valori aziendali fondamentali di CodicePlastico vengono annoverati l'aggiornamento costante del personale e il miglioramento delle \textit{skill} di ogni sviluppatore, anche attraverso pratiche agili come \textit{pair programming}, \textit{randori} e progetti personali. Lo sviluppo di una applicazione con l'ausilio di tecnologie relativamente recenti e innovative nel panorama IT, � in completo accordo con la \textit{vision} aziendale.
\end{description}
	\clearpage

\chapter{Metodologie di sviluppo}\label{Metodologie_di_sviluppo}
	\section{User stories}\label{User_stories}
	\clearpage
	\section{Test e behavior driven development}\label{Test_e_behavior_driven_development}
	\clearpage

\chapter{Tecnologie utilizzate}\label{Tecnologie_utilizzate}
	\section{Node.js}\label{Nodejs}
		% motivo della scelta
		% peculiarit�
		\subsection{La storia}\label{La_storia}
		\clearpage
		\subsection{NPM e moduli}\label{NPM_e_moduli}
		% cos'� npm, cos'� un modulo
		% gestione delle dipendenze
		% alcuni dei moduli che abbiamo usato (i pi� significativi)
		\clearpage
	\section{RabbitMQ}\label{RabbitMQ}
	% motivo della scelta (disaccoppiamento, vantaggi)
	% configurazione utilizzata per scalare (broadcast+loadbalancing)
	\clearpage
	\section{MongoDB}\label{MongoDB}
	% motivo della scelta
		\subsection{Fronteggiare le richieste}\label{Fronteggiare_le_richieste}
		\clearpage
	\section{AngularJS e altre tecnologie di frontend}\label{AngularJS_e_altre_tecnologie_di_frontend}
	% motivo della scelta
		\subsection{Gestione delle dipendenze}\label{Gestione_delle_dipendenze}
		\clearpage
	\section{Strumenti per il deploy}\label{Strumenti_per_il_deploy}
	\clearpage

\chapter{Mole.io}\label{Moleio}
	\section{Architettura del sistema}\label{Architettura_del_sistema}
		\subsection{CQRS ed estensibilit�}\label{CQRS_ed_estensibilita}
		\clearpage
		\subsection{mole}\label{mole}
			\subsubsection{I denormalizzatori}\label{I_denormalizzatori}
			\clearpage
		\subsection{mole-suit}\label{mole-suit}
			\subsubsection{I plugin e gli widget}\label{I_plugin_e_gli_widget}
			\clearpage
	\section{Autenticazione degli utenti}\label{Autenticazione_degli_utenti}
	\clearpage
	\section{Scalabilit� e affidabilit�}\label{Scalabilita_e_affidabilita}
	\clearpage
	\section{Problematiche di sviluppo}\label{Problematiche_di_sviluppo}
	\clearpage

\chapter{Configurazioni e benchmark}\label{Configurazioni_e_benchmark}
problemi con i benchmark
- dipendi dalla rete su cui sei
- nostro client fatto con node - perch� non l'abbiamo usato
- 2 parole sul tool che abbiamo usato
- come sono stati fatti i benchmark
- specifiche del sistema VM, ram, hdd, ...
- risultati ottenuti


non ha senso testare mole (server) per la parte rabbit (spiegare)
il vero collo di bottiglia � il db
abbiamo testato con configurazioni di db differenti

% parlare dei casini con le chiavi di shard



\clearpage

\chapter*{Conclusioni e sviluppi futuri\markboth{}{Conclusioni e sviluppi futuri}}\label{Conclusioni_e_sviluppi_futuri}
\addcontentsline{toc}{chapter}{Conclusioni e sviluppi futuri}
testere il sistema "in piccolo" � difficile, provare test sul sistema deployato

%TODO
% workflow con webhooks
% notifiche push (si potrebbe usare un denormalizer?)
% app mobile per riceverle
% https
% marcatura degli errori come "gestito" e conseguente archiviazione
% integrazione con sistemi di tracking (?)
% aggiungere pi� mole-contacts
% aggiungere un sistema di query
% mole-contacts -> caching degli errori in locale quando non c'� rete (buffering)

% migliorare la gestione dgli widget e plugin per l'aggiunta a caldo

% possibili nuovi denormalizzatori

% bench e vm sulla stessa macchina
% vm con cpu pin
% ramp up

\clearpage

\bibliography{Bibliografia}
\bibliographystyle{unsrt}
\addcontentsline{toc}{chapter}{Bibliografia}

\end{document}



%
%il problema
%- problema del logging:
%  - tante applicazioni
%  - diversi clienti
%  - report errori
%
%stato dell'arte: i software di logging centralizzato
%- aircoso
%- fluentd + elastic search + kibana
%perch\'{e} lo stato dell'arte non ci piace
%- non permettono di avere "modelli/template di messaggi" flessibili
%- contesti differenti
%
%il nostro logger 
%- perch\'{e} abbiamo scelto nodejs
%- cos'� nodejs (storia, peculiarit�, async, pesca da presentazione lucio)
%- come funziona il logger (struttura e schema delle comunicazioni di rete)
%- problematiche incontrate durante l'implementazione
%- si auto-adatta al tipo di messaggio in arrivo
%- denormalizzatori
%- plugin/widget
%- il clAud
%- architettura del sistema
%- estensibilit�
%- come scrivere plugin e widget
%
%problematiche "accademiche"
%- sicurezza (oauth)
%- scalabilit� (+ server con load balancer? cloud? rabbit, disaccoppiamento dei servizi)
%- relayability (+ server ridondati) 
%
%
%
%Scaletta:
%
%Metodologie:
% - CQRS
% - TDD
% - user stories
%
%Tecnologie utilizzate:
%- NodeJS
%  - npm (gestione dipendenze)
%  - passport (oauth)
%  - mongoskin
%  - mongoose
%  - underscore
%  - express
%  - faker (fixtures)
%  - grunt (automation)
%  - mocha (testing)
%  - should (testing)
%  - supertest (testing)
%  - require-all (plugin)
%  - cors (e problemi correlati)
%
%- rabbitmq
%
%- angularjs
%  - leflet (mappe)
%  - directive
%  - controller
%  - views
%  - services
%  - filters
%  - karma (testing)
%  - d3 (grafici)
%  - jquery
%  - bootstrap (ui)
%  - bower (gestione dipendenze)
%
%- mongodb
%  - sharding
%  - replicaset
%
%- dokku
%  
%- git(?)





